\section{Riemannian Manifolds}

\subsection{Riemannian manifolds}

A symmetric $2$-tensor $\alpha\in\sym^2(V^*)$ is {\emphcolor positive definite} if $\alpha(v,v)>0$ for all $v\neq0$.
A symmetric covariant $2$-tensor field is positive definite if it is positive definite at every point.

\bdefn

    Let $M$ be a smooth manifold, a {\emphcolor Riemannian metric} on $M$ is a smooth symmetric covariant $2$-tensor field on $M$.
    A {\emphcolor Riemannian manifold} is a smooth manifold $M$ equipped with a Riemannian metric $g$.

\edefn

Since $g_p$ is an positive semidefinite and symmetric, it forms an inner product on $T_pM$.
We denote this by $\iprod{\bul,\bul}_g$.

\bnote

    Equivalently, we can characterize a Riemannian metric to be an association to each $p\in M$ an inner product $\iprod{\bul,\bul}_g$ on $T_pM$
    which is smooth in the sense that $\iprod{X,Y}_g$ is in $C^\infty(M)$ for every $X,Y\in\X(M)$.

\enote

In any smooth coordinate chart $(x^i)$, a Riemannian metric has components
$$ g = g_{ij}dx^i\otimes dx^j . $$
Note then that by symmetry, $g_{ij}=g_{ji}$ and so
$$ g = \frac12(g_{ij}dx^i\otimes dx^j + g_{ij}dx^i\otimes dx^j) = \frac12(g_{ij}dx^i\otimes dx^j + g_{ji}dx^i\otimes dx^j) =
\frac12(g_{ij}dx^i\otimes dx^j + g_{ij}dx^j\otimes dx^i) = g_{ij}dx^idx^j , $$
where $dx^idx^j$ is the symmetric product.

\bexam

    The {\emphcolor Euclidean metric} on $\bR^n$ is $\bar g$ given by gloabl coordinates
    $$ \bar g = \delta_{ij}dx^idx^j . $$
    Writing $\alpha^2$ for the symmetric product of $\alpha$ with itself, we have that
    $$ \bar g = (dx^1)^2 + \cdots + (dx^n)^2 . $$
    Then for $v,w\in T_p\bR^n$ we have
    $$ \bar g_p(v,w) = \delta_{ij}v^iw^j = \sum_i v^iw^i = v\cdot w, $$
    so that the Euclidean metric is the Euclidean inner product on $T_p\bR^n$.

\eexam

\bexam

    If $(M_1,g^1)$ and $(M_2,g^2)$ are Riemannian manifolds, we can define a Riemannian metric $g=g^1\oplus g^2$ on $M_1\times M_2$ by
    $$ g((v_1,w_1),(v_2,w_2)) = g^1(v_1,v_2) + g^2(w_1,w_2) , $$
    for $(v_i,w_i)\in T_pM_1\oplus T_qM_2\cong T_{(p,q)}(M_1\times M_2)$.
    This is the sum of two smooth tensor fields, and thus smooth itself, and symmetric.
    It is positive definite at each point because $g^1,g^2$ are.

    Given local coordinates $(x^i)$ for $M_1$ and $(y^j)$ for $M_2$, if $g^1$ has components $g^1_{ij}dx^idx^j$ and $g^2$ has components $g^2_{ij}dy^idy^j$, then
    $g$'s coordinate representation is
    $$ (g_{ij}) = \pmatrix{g^1_{ij}\cr&g^2_{ij}} . $$

\eexam

\bprop

    Every smooth manifold admits a Riemannian metric.

\eprop

\bproof

    Let $(U_\alpha,\phi_\alpha)$ cover $M$ by countably many coordinate charts.
    Then $\phi_\alpha\colon U_\alpha\to\bR^m$ is smooth and thus we can pull back the Euclidean metric to $U_\alpha$ by $g_\alpha=\phi_\alpha^*\bar g$.
    Its coordinate representation is $\delta_{ij}dx^idx^j$ relative to $\phi_\alpha$.
    Now let $\set{\psi_\alpha}$ be a partition of unity subordinate to $U_\alpha$, and define
    $$ g = \sum_\alpha\psi_\alpha g_\alpha . $$
    This is smooth, by local finiteness.
    It is also obviously symmetric, so only positivity needs to be checked.
    All the functions in question are nonnegative, so we need to show that $g(v,v)\neq0$ when $v\neq0$.
    Indeed,
    $$ g(v,v) = \sum_\alpha\psi_\alpha(p)g_\alpha(v,v) , $$
    and if $0\neq v\in T_pM$ then some $\psi_\alpha(p)$ must be positive (since they sum to $1$), and thus $g(v,v)\geq\psi_\alpha(p)g_\alpha(v,v)>0$.
    \qed

\eproof

Note that such a choice of Riemannian metric is not natural.

\bdefn

    If $(M,g)$ is a Riemannian metric, then the {\emphcolor norm} of $v\in T_pM$ is defined to be $\norm v=\iprod{v,v}_g^{1/2}=g_p(v,v)^{1/2}$.
    $v,w\in T_pM$ are {\emphcolor orthogonal} if $\iprod{v,w}_g=0$.

\edefn

\bprop

    Let $(M,g)$ be a Riemannian manifold, and $(X_j)$ a smooth local frame for $M$ over an open subset $U\subseteq M$.
    Then there is a smooth local frame $(E_j)$ over $U$ which is orthonormal at every point, and $\lspan(E_1|_p,\dots,E_j|_p)=\lspan(X_1|_p,\dots,X_j|_p)$
    for all $j=1,\dots,n$.

    In particular for every $p\in M$ there is an orthonormal frame in a neighborhood around $p$.

\eprop

\bproof

    Apply Gram-Schmidt, and note that the resulting fields are smooth.
    \qed

\eproof

\bdefn

    If $F\colon M\to N$ is smooth, and $g$ is a Riemannian metric on $N$, then the pullback $F^*g$ is a smooth $2$-tensor field on $M$.
    If it is positive definite, it is a Riemannian metric on $M$, called the {\emphcolor pullback metric}.

\edefn


\bprop

    The pullback $F^*g$ is Riemannian iff $F$ is a smooth immersion.

\eprop

\bproof

    The pullback is defined by
    $$ F^*g_p(v,w) = g(dF_pv,dF_pw) , $$
    and in general $\iprod{F\bul,F\bul}$ is an inner product only when $F$ is injective.
    \qed

\eproof

\bdefn

    If $(M,g)$ and $(N,\tilde g)$ are Riemannian manifolds, a smooth map $F\colon M\to N$ is a {\emphcolor (Riemannian) isometry} if it is a diffeomorphism
    which satisfies $F^*\tilde g=g$.
    A {\emphcolor local isometry} is one where for every $p\in M$ there is a neighborhood $p\in U$ such that $F|_U\colon U\to N$ is an isometry onto an open subset of $N$.
    Equivalently, if it satisfies $F^*\tilde g=g$ and $F$ is a local diffeomorphism.

    If there exists an isometry between Riemannian manifolds, we say they are {\emphcolor isometric}.
    A Riemmanian $n$-manifold is a {\emphcolor flat Riemannian manifold} (and the metric is a {\emphcolor flat metric}) if it is locally isometric to $(\bR^n,\bar g)$.

\edefn

\bprop

    Let $(M,g)$ be a Riemannian manifold.
    The following are equivalent:
    \benum
        \item $g$ is flat;
        \item Each point of $M$ is contained in the domain of a smooth coordinate chart in which $g$ has the coordinate representation $g=\delta_{ij}dx^idx^j$;
        \item Each point of $M$ is contained in the domain of a smooth coordinate chart in which the coordinate frame is orthonormal;
        \item Each point of $M$ is contained in the domain of a commuting orthnormal frame.
    \eenum

\eprop

\bproof

    Clearly (2) and (3) are equivalent, and if $g$ is flat then every point lies in an open subset $U\subseteq M$ which is
    isometric to an open subset of $\bR^n$.
    The isometry forms a local chart $(U,(x^i))$ where $\ipdv{x^i}$ are orthonormal.
    And conversely if $(U,(x^i))$ is a local chart where $\ipdv{x^i}$ are orthonormal, then the chart is an isometry as it
    maps the local coordinate frame $\ipdv{x^i}$ to the Euclidean coordinate frame which is orthonormal.
    It is sufficient for a map to preserve a single orthonormal frame for it to be an isometry.

    Clearly (3) implies (4) as coordinate frames are commuting.
    Now (4) implies (3) since by the canonical form for commuting frames, if $(E_i)$ is a commuting local orthonormal frame
    then there is a coordinate frame $(U,(x^i))$ such that $E_i=\ipdv{x^i}$.
    \qed

\eproof

\bdefn

    If $S\subseteq M$ is an immersed submanifold with $(M,g)$ Riemannian, then $S$ inherits the pullback metric $\iota^*g$ where $\iota\colon S\inj M$ is the inclusion.
    This is called the {\emphcolor induced metric}.

\edefn

The induced metric is indeed a metric since $\iota$ is by definition an immersion.
Under this metric $S$ is called a {\emphcolor Riemannian submanifold}.

\bexam

    The induced metric $\iota^*g$ induced on $S^n$ by the inclusion $S^n\inj\bR^{n+1}$ is called the {\emphcolor round} or {\emphcolor standard metric} on the sphere.

\eexam

\bdefn

    Let $(M,g)$ be a Riemannian $n$-manifold, and $S\subseteq M$ a $k$-dimensional submanifold.
    We say that for $p\in S$, a vector $v\in T_pM$ is {\emphcolor normal to $S$} if it is orthogonal to $T_pS$ with respect to $\iprod{\bul,\bul}_g$.
    The {\emphcolor normal space to $S$} at $p$ is the subspace $N_pS\subseteq T_pM$ of all vectors orthogonal to $T_pS$ (i.e.\ $N_pS=(T_pS)^\perp$).

    Let $NS=\coprod_{p\in S}N_pS$, and $\pi_{NS}\colon NS\to S$ to be the obvious projection.

\edefn

\bprop

    $NS\to S$ forms a smooth $(n-k)$-rank vector bundle over $S$, with a unique smooth structure where the linear structures on the fibers are given by $N_pS$.
    For each $p\in S$ there is a smooth frame for $NS$ on a neighborhood of $p$ that is orthonormal with respect to $g$.

\eprop

\bproof

    Take a coordinate frame $(X_i)$ for $T_pS$, and then we can extend it and apply Gram-Schmidt to get a coordinate frame $(X_i,Y_j)$ where $Y_j$ are in $N_pS$.
    These form a basis, and thus satisfy the local frame criterion for subbundles.
    \qed

\eproof

\subsection{The Riemannian distance function}

\bdefn

    Suppose $(M,g)$ is a Riemannian manifold, and $\gamma\colon[a,b]\to M$ a piecewise smooth curve segment.
    $\gamma$'s {\emphcolor length} is
    $$ L_g(\gamma) = \int_a^b\norm{\gamma'(t)}_g\,dt . $$
    This is well-defined since $\gamma'(t)$ is defined at all but finitely many values of $t$.

\edefn

Note that if $F\colon M\to N$ is a local isometry, then since $F^*g_N=g_M$ we have that $F$ preserves lengths: if $\gamma\colon[a,b]\to M$ is piecewise smooth,
$$ L_{g_N}(F\circ\gamma) = \int_a^b\norm{(F\circ\gamma)'(t)} = \int_a^b\norm{dF_{\gamma(t)}\circ\gamma'(t)} = \int_a^b\norm{\gamma'(t)} = L_{g_M}(\gamma). $$

\bdefn

    A {\emphcolor reparametrization} of a curve $\gamma\colon[a,b]\to M$ is a curve of the form $\tilde\gamma=\gamma\circ\phi$ where $\phi\colon[c,d]\to[a,b]$ is
    a diffeomorphism.

\edefn

\bprop

    Curve length is invariant under reparametrization.

\eprop

\bproof

    Let $\phi\colon[c,d]\to[a,b]$ be a diffeomorphism and $\gamma\colon[a,b]\to M$ a piecewise smooth curve.
    Since $\phi$ is a diffeomorphism, either $\phi'>0$ or $\phi'<0$ everywhere.
    Assume $\phi'$ is always positive, then
    $$ L_g(\gamma\circ\phi) = \int_c^d\norm{(\gamma\circ\phi)'(t)} = \int_c^d\norm{\phi'(t)\gamma'(\phi(t))} = \int_c^d\norm{\gamma'(\phi(t))}\phi'(t) = \int_a^b\norm{\gamma'(t)} , $$
    where the final equality is due to change of variables.
    If $\phi'<0$ then we'd have to multiply the second to last term by $-1$, then change of variables swaps the sign again.
    \qed

\eproof

\bdefn

    Let $(M,g)$ be a connected Riemannian manifold.
    For $p,q\in M$ define their {\emphcolor Riemannian distance}, denoted $d_g(p,q)$, to be the infimum over $L_g(\gamma)$ over all piecewise smooth curves $\gamma$
    from $p$ to $q$.

\edefn

\bexam

    Consider $\bR^n$ with the Euclidean metric $g$.
    Then $d_g(p,q)=\norm{p-q}$, i.e.\ the Riemannian distance is the same as Euclidean distance.

    Indeed, $\gamma\colon[0,1]\to\bR^n$ with $\gamma(t)=tp+(1-t)q$ has $\gamma'(t)=p-q$ and thus
    $$ L_g(\gamma) = \int_0^1\norm{p-q} = \norm{p-q} . $$
    So $d_g(p,q)\leq\norm{p-q}$.

    Now if $\gamma$ is any other piecewise smooth curve, we need to show that $L_g(\gamma)\geq\norm{p-q}$.
    We have that
    $$ L_g\gamma = \int_a^b\norm{\gamma'(t)}\,dt \geq \norm{\int_a^b\gamma'(t)} = \norm{\gamma(b) - \gamma(a)} = \norm{p-q} , $$
    as desired.

\eexam

\blemm

    Let $g$ be a Riemannian metric on an open subset $U\subseteq\bR^n$.
    Given a compact subset $K\subseteq U$, there exist positive constants $c,C$ such that for all $x\in K$ and $v\in T_x\bR^n$,
    $$ c\norm v_{\bar g} \leq \norm v_g \leq C\norm v_{\bar g} . $$

\elemm

\bproof

    Define $L\subseteq T\bR^n$ by
    $$ L = \set{(x,v)\in T\bR^n}[x\in K,\norm v_{\bar g}=1] . $$
    Under the canonical identification of $T\bR^n$ with $\bR^n\times\bR^n$, this is just $L=K\times S^{n-1}$, which is compact.
    Since $\norm v_g$ is continuous and strictly positive on $L$, there are positive constants $c,C$ such that $c\leq\norm v_g\leq C$ ($\norm v_g$
    obtains its maximum and minimum on a compact set) for all $(x,v)\in L$.
    If $x\in K$ and $v\in T_x\bR^n$ is nonzero, let $\lambda=\norm v_{\bar g}$ so that $(x,\lambda^{-1}v)\in L$ so
    $$ \norm v_g = \lambda\norm{\lambda^{-1}v}_g \leq \lambda C = C\norm v_{\bar g}, $$
    and similar for $c$.
    \qed

\eproof

\bthrm

    The Riemannian distance function on a connected Riemannian manifold $(M,g)$ is a metric in the usual sense.
    It induces the same topology on $M$ as its manifold topology.

\ethrm

\bproof

    All of the conditions for $d_g$ being a metric are clear, other than positivity.
    Note that the triangle inequality is because if $\gamma$ connects $a$ to $b$, and $\tilde\gamma$ connects $b$ to $c$, then we can concatenate the curves
    to get a curve from $a$ to $c$ whose length is the sum of the curves.

    Now we need to show that $d_g(p,q)\neq0$ when $p\neq q$.
    Let $U$ be a chart containing $p$ but not $q$.
    We use the coordinate map to identify $U$ with a an open subset of $\bR^n$, and let $\bar g$ be the Euclidean metric in these coordinates.
    Let $V$ be a regular coordinate ball of radius $\epsilon$ centered at $p$ such that $\cl V\subseteq U$.
    Then the above lemma shows that there are constants $c,C>0$ such that $c\norm v_{\bar g}\leq\norm v_g\leq C\norm v_{\bar g}$ for all $v\in T_xM$ when $x\in\cl V$.
    Then for any smooth curve $\gamma$ lying in $\cl V$, $cL_{\bar g}(\gamma)\leq L_g(\gamma)\leq CL_{\bar g}(\gamma)$.

    Now suppose $\gamma\colon[a,b]\to M$ is piecewise smooth from $p$ to $q$.
    Let $t_0$ be the infimum of all $t$ such that $\gamma(t)\notin\cl V$.
    By continuity, $\gamma(t_0)\in\partial V$ and so $\gamma(t)\in\cl V$ for $a\leq t\leq t_0$.
    Thus
    $$ L_g(\gamma) \geq L_g(\gamma|_{[a,t_0]}) \geq cL_{\bar g}(\gamma|_{[a,t_0]}) \geq cd_{\bar g}(p,\gamma(t_0)) = c\epsilon . $$
    The final inequality is because $d_{\bar g}$ agrees with normal Euclidean distance, and $\norm{p-\gamma(t_0)}=\epsilon$ since $\gamma(t_0)\in\partial V$.
    Thus we conclude that $d_g(p,q)\geq c\epsilon>0$ as desired.

    Finally to show that the two topologies coincide, we show that open sets in one are open in the other and vice versa.
    Suppose that $U\subseteq M$ is open, and let $p\in U$, and let $V$ be a regular coordinate ball of radius $\epsilon$ around $p$ such that $\cl V\subseteq U$.
    We have that $d_g(p,q)\geq c\epsilon$ when $q\notin\cl V$, and thus if $d_g(p,q)<c\epsilon$ then $q\in\cl V\subseteq U$.
    So the open ball of radius $c\epsilon$ centered around $p$ is contained in $U$, since $p$ was arbitrary this means $U$ is open.

    Conversely if $W$ is open in the metric topology, let $p\in W$.
    Let $V$ be a regular coordinate ball around $p$ of radius $r$, and let $\bar g$ be the Euclidean metric on $\cl V$.
    Then we have positive constants $c,C>0$ bounding the metric of $d_g$ by $d_{\bar g}$ as before.
    Let $\epsilon<r$ be small enough such that the metric ball around $p$ of radius $C\epsilon$ is contained in $W$, and let $V_\epsilon$ be the set of points of $\cl V$
    whose Euclidean distance from $p$ is less than $\epsilon$.
    If $q\in V_\epsilon$ let $\gamma$ be the straight line segment from $p$ to $q$, then
    $$ d_g(p,q) \leq L_g(\gamma) \leq CL_{\bar g}(\gamma) < C\epsilon . $$
    Thus $V_\epsilon$ is contained in the metric ball of radius $C\epsilon$ around $p$, and thus in $W$.
    $V_\epsilon$ is open in the manifold topology (since it is an open set under the coordinate chart of $\cl V$), and $p\in V_\epsilon$, we see that $W$ is open in the
    manifold topology.
    \qed

\eproof

\bdefn

    A Riemannian manifold $(M,g)$ is {\emphcolor complete}, and $g$ is a {\emphcolor complete Riemannian metric}, if $(M,d_g)$ is a complete metric space.
    A subset $B\subseteq M$ is {\emphcolor bounded} if it is bounded in the metric space.

\edefn

\bcoro

    Every smooth manifold is metrizable.

\ecoro

\bproof

    If $M$ is a manifold without boundary, then choose a Riemannian metric $g$ on $M$.
    If $M$ is connected, then the above proposition shows that it is metrizable by $d_g$.
    Otherwise, let $\set{M_i}$ be the components of $M$, each of which is metrizable by $d_g|_{M_i}$.
    The disjoint union of metrizable spaces is metrizable: let $p_i\in M_i$ and define
    $$ d_g(x,y) = d_g(x,p_i) + 1 + d_g(p_j,y), $$
    when $x,y$ are not in the same component.

    If $M$ has boundary, it can be embedded into a manifold without boundary (its double), which is metrizable.
    Subspaces of metrizable spaces are metrizable.
    \qed

\eproof

\subsection{The musical isomorphisms}

If $V$ is a finite-dimensional inner product space, there is a natural isomorphism $V\to V^*$ given by $v\mapsto\iprod{v,\bul}$.
This is clearly linear and injective, since if $\iprod{v,\bul}=0$ then $\iprod{v,v}=0$ so $v=0$.
We can also compute its inverse: let $e_i$ be an orthonormal basis of $V$, then for $\phi\in V^*$ define $v=\sum_i\phi(e_i)e_i$.

Thus have a bundle isomorphism
$$ \hat g\colon TM\to T^*M,\qquad \hat g(v)(w) = g(v,w) . $$
This is a bundle isomorphism since it corresponds to the morphism of sections $\hat g\colon\X(M)\to\X^*(M)$, $\hat g(X)(Y)=g(X,Y)$.
Indeed, $\hat g(X)$ is a smooth covector field because its action on vector fields is smooth.
And clearly $\hat g$ is $C^\infty(M)$-linear.
Furthermore, $\hat g$ is bijective on the fibers, so that $\hat g$ is a bijective bundle morphism over $M$ and thus an isomorphism.

If $(E_i)$ is a local frame with dual coframe $(\epsilon^i)$, let $g=g_{ij}\epsilon^i\epsilon^j$ be the representation of the metric with respect to this frame.
Suppose $X\in\X(M)$ is a vector field with representation $X=X^iE_i$, then
$$ \hat g(X)(E_j) = g(X^iE_i,E_j) = X^ig_{ij} $$
so that
$$ \hat g(X) = (X^ig_{ij})\epsilon^j . $$
It is customary to write the components of $\hat g(X)$ as $X_j$, so that $X_j=X^ig_{ij}$.
Due to this notational convention, the action of $\hat g$ on vector fields is called ``lowering an index'', and we write $X^\flat$ for $\hat g(X)$.
The symbol $\flat$ is a musical symbol called the {\emphcolor flat} which denotes the lowering of a tone, and thus $X^\flat$ is called $X$-flat.

The inverse $\hat g^{-1}$ acts on covector fields $\omega=\omega_j\epsilon^j$ and maps them to vector fields
$$ \hat g^{-1}(\omega) = \omega^iE_i,\qquad \omega^i = g^{ij}\omega_j , $$
where $(g^{ij})$ is the inverse matrix of $(g_{ij})$.
This is called ``raising an index'', and we write $\omega^\sharp$ for $\hat g^{-1}(\omega)$; borrowing the musical symbol for raising a tone.

Thus we have bijective correspondences

\medskip
\centerline{\vbox{\openup2\jot\halign{\hfil$#$&\hfil$#$\hfil&$#$,\hfil\tabskip=1cm&$#$\hfil\cr
    \X(M)&{}\xtof{\quad\bul^\flat\quad}[\quad\bul^\sharp\quad]{}&\X^*(M)\cr
    X^iE_i&\xmapsto{\qquad}&X_i\epsilon^i&X_i=g_{ij}X^j\cr
    \omega^iE_i&\xmapsfrom{\qquad}&\omega_i\epsilon^i&\omega^i=g^{ij}\omega_j.\cr
}}}
\medskip

These correspondences are called the {\emphcolor musical isomorphisms}.
Note that we can raise and lower the indices of non-smooth vector and covector fields in the same way.

\subsubsection*{The gradient}

One application of the musical isomorphisms is to recover the notion of the gradient of a real-valued function.
Recall that if $f\colon\bR^n\to\bR$ is a smooth function, its gradient $\grad f$ is the vector of its partial derivatives.
It is the vector with the unique property that
$$ \grad f|_p\cdot v = D_vf , $$
where $D_vf$ is the directional derivative of $f$ in the direction of $v$.

In the case of Riemannian manifolds, $D_vf$ is replaced with $df(v)$ for $v\in T_pM$ so that this becomes
$$ \iprod{\grad f,v}_p = df_p(v) , $$
which is to say that $\grad f=(df)^\sharp$.
Thus in any local frame $(E_i)$ with dual frame $(\epsilon^i)$, we have that $df=(E_if)\epsilon^i$ and so
$$ \grad f = (g^{ij}E_if)E_j . $$
So if $(E_i)$ is an orthonormal frame, the components of $\grad f$ are the same as the components of $df$; but this is not true in general.
In the case of $\bR^n$ we see that $\grad f=\ipdvof f{x_i}\ipdv{x_i}$.

\subsubsection*{The musical isomorphisms on tensors}

Now we can generalize the musical isomorphisms to any tensor.
If $F$ is a $(k,\ell)$-tensor, then it corresponds to a map
$$ F\colon\X^*(M)^{\otimes k}\otimes\X(M)^{\otimes\ell} \to C^\infty(M) , $$
and we can use the musical isomorphisms on any index $i\in\set{1,\dots,k+\ell}$ to change a $\X(M)$ to a $\X^*(M)$ or vice versa.
That is, if $i$ is a covariant index, meaning its input is a vector, then we obtain $F^\sharp$ by making that index contravariant.
Explicitly, $F^\sharp$ is a $(k+1,\ell-1)$-tensor obtained by
$$ F^\sharp(\alpha_1,\dots,\alpha_{k+\ell}) = F(\alpha_1,\dots,\alpha_i^\sharp,\dots,\alpha_{k+\ell}) . $$
Similarly if $i$ is a contravariant index, then $F^\flat$ is obtained by making that index covariant:
$$ F^\flat(\alpha_1,\dots,\alpha_{k+\ell}) = F(\alpha_1,\dots,\alpha_i^\flat,\dots,\alpha_{k+\ell}) . $$
In this way $F^\flat$ is a $(k-1,\ell+1)$-tensor.

Note that $X\in\X(M)$ is a $(1,0)$-tensor and this definition makes $X^\flat$ a $(0,1)$-tensor (covector field) by
$$ X^\flat(Y) = X(Y^\flat) = Y^\flat(X) = \hat g(Y)(X) = g(X,Y) $$
which is the same as our previous definition of $X^\flat$.

In terms of a local frame, $F^\sharp$ is obtained by multiplying the components of $F$ with $g^{pq}$ where we set either $p$ or $q$ to be the index
of $F$ which we raise.
Similar for $F^\flat$.
For example, if we have a mixed $3$-tensor
$$ A = {{A_i}^j}_k\epsilon^i\otimes E_j\otimes\epsilon^k , $$
and we want to lower the middle index to obtain a covariant $3$-tensor $A^\flat$, we get
$$ A^\flat = A_{ijk}\epsilon^i\otimes\epsilon^j\otimes\epsilon^k,\qquad A_{ijk} = g_{j\ell}{{A_i}^\ell}_k . $$

While just writing $F^\sharp,F^\flat$ is ambiguous, as we have not specified which index to raise/lower, we will often omit the index from our notation and instead
specify it in words when necessary.

\subsubsection*{The trace of a tensor}

We turn to another application of the musical isomorphisms.
Recall that $T^{(1,1)}V=V\otimes V^*\cong\hom(V,V)$, so that we can unambiguously take the trace of a $(1,1)$-tensor.
Explicitly, if $F=F^i_jE_i\otimes\epsilon^j$ is a $(1,1)$-tensor, then its trace is the sum
$$ \tr(F) = F^i_i . $$
We can extend this to mixed tensors: if $F$ is a mixed $(k+1,\ell+1)$-tensor, for any choice of a covariant index and contravariant index we can form $F$'s trace.
We can assume that the choice of indices are the last ones; then for covectors $\omega^1,\dots,\omega^k$ and vectors $v_1,\dots,v_\ell$, we have a $(1,1)$-tensor
$$ F(\omega^1,\dots,\omega^k,\bul,v_1,\dots,v_\ell,\bul) \in T^{(1,1)}V . $$
Thus we can define $\tr(F)(\bar\omega,\bar v)$ to be the trace of this tensor.
This defines a $(k,\ell)$-tensor.
If $F$ has components $F^{I,i}_{J,j}$, then $\tr(F)$'s components are the sums $F^{I,m}_{J,m}$.

So $\tr$ defines a bundle morphism
$$ \tr\colon T^{(k+1,\ell+1)}TM \to T^{(k,\ell)}TM . $$

Thus far we have not had to use the musical morphisms.
But if we have a covariant $k$-tensor $h$ with $k\geq2$, then $h^\sharp$ is a $(1,k-1)$-tensor for any choice of index to raise.
We can then take its trace, resulting in a covariant $(k-2)$-tensor,
$$ \tr_gh = \tr(h^\sharp) . $$
If $h$'s components are $h=h_I\epsilon^I$, then $h^\sharp=g^{ij}h_{i_1,\dots,i_{k-1},i}E_j\otimes\epsilon^{i_1}\otimes\cdots\otimes\epsilon^{i_{k-1}}$.
Thus $\tr_gh$ has components
$$ (\tr_gh)_{i_1,\dots,i_{k-2}} = (h^\sharp)_{i_1,\dots,i_{k-2},j}^j = g^{ij}h_{i_1,\dots,i_{k-2},j} . $$

For example if $h$ is a $2$-covector field, so $h=h_{i,j}\epsilon^i\otimes\epsilon^j$ then $h^\sharp=g^{i\ell}h_{\ell,j}E_i\otimes\epsilon^j$ so that
$$ \tr_gh = g^{i\ell}h_{\ell i} $$
is a smooth function.

Note that we can take the trace of all sorts of tensors.
For example, a tensor valued in another manifold.
We will revisit this when necessary, e.g.\ with mean curvature.

\subsubsection*{Metrics on tensors}

We can also use the musical isomorphisms to define the inner product of tensors.
If $M$ is Riemannian, then we can define an inner product on the cotangent space $T_p^*M$ for $p\in M$ by
$$ \iprod{\omega,\eta}_g = \iprod{\omega^\sharp,\eta^\sharp}_g . $$
Note that applying this inner product on smooth covector fields results in a smooth real-valued function (by smoothness of $g$).

In terms of a local frame $(E_i)$ with dual coframe $(\epsilon^j)$, if $\omega=\omega_i\epsilon^i$ and $\eta=\eta_j\epsilon^j$, then
$$ \iprod{\omega,\eta}_g = \iprod{g^{ik}\omega_k E_i,g^{j\ell}\eta_\ell E_j} = g^{ik}g^{j\ell}\omega_k\eta_\ell g_{ij} = g^{ij}\omega_i\eta_j $$
where the last equality is because $g^{ik},g_{ij}$ are inverses.
Using our convention for raising indices, where $\omega^i=g^{ij}\omega_j$, we see that
$$ \iprod{\omega,\eta} = \omega^i\eta_i = \omega_i\eta^i . $$

Note that it is not the case in general that $(\epsilon^i)^\sharp=E_i$.
This is true iff $(E_i)$, equivalently $(\epsilon^i)$, is orthonormal.
Indeed, $(\epsilon^i)^\sharp=g^{ij}E_j$ is equal to $E_i$ iff $g^{ij}$ is the identity, iff $g_{ij}$ is, which means $(E_i)$ is orthonormal.

We can naturally define a metric on a vector bundle $E\to M$ to be a metric on each fiber $E_p$ which varies smoothly as $p$ does, in the sense that for any local
sections $\sigma,\tau$, $\iprod{\sigma,\tau}$ is a smooth function.
Equivalently, this is just a section of $\Gamma(E^*\otimes E^*)$ which is symmetric and positive-definite on each fiber.

If $E_1,E_2\to M$ are vector bundles over $M$ with smooth metrics, we can define a smooth metric on $E_1\otimes E_2$ by
$$ \iprod{\alpha\otimes\beta,\alpha'\otimes\beta'} = \iprod{\alpha,\alpha'}\cdot\iprod{\beta,\beta'},\qquad \alpha,\alpha'\in\Gamma(E_1),\quad \beta,\beta'\in\Gamma(E_2) . $$
This can be viewed as being obtained via the isomorphism
$$ \Gamma((E_1\otimes E_2)^*\otimes(E_1\otimes E_2)^*) \cong \Gamma(E_1^*\otimes E_1^*)\otimes\Gamma(E_2^*\otimes E_2^*) . $$

Since we have metrics on $TM,T^*M$, by this result we can define a unique inner product on the tensor bundles $T^{(k,\ell)}M$.
It is given uniquely by
$$ \iprod{\alpha_1\otimes\cdots\otimes\alpha_{k+\ell},\beta_1\otimes\cdots\otimes\beta_{k+\ell}} = \iprod{\alpha_1,\beta_1}\cdots\iprod{\alpha_{k+\ell},\beta_{k+\ell}} $$
where $\alpha_i,\beta_i$ are vector or covector fields, where appropriate.
We see that if $F,G$ are $(k,\ell)$-tensors with components $F^I_J,G^I_J$, then
$$ \iprod{F,G} = g_{i_1,r_1}\cdots g_{i_k,r_k}g^{j_1,s_1}\cdots g^{j_\ell,s_\ell}F^{i_1,\dots,i_k}_{j_1,\dots,j_\ell}G^{r_1,\dots,r_k}_{s_1,\dots,s_\ell} . $$

